% Capítulo 4: Conclusões e Trabalhos Futuros
\chapter{Conclusões}\label{chp:conclusoes}

Esta dissertação apresentou uma metodologia sistemática para otimização computacional de processamento hiperespectral embarcado, focando na análise comparativa de trade-offs entre arquiteturas heterogêneas FPGA e GPU aplicadas à agricultura de precisão. Partindo da identificação de uma lacuna metodológica significativa na literatura, onde apenas 4\% dos 25 trabalhos revisados (2011-2024) abordam sistemas completos integrados para aplicações agrícolas, desenvolveu-se um framework unificado estruturado em oito componentes principais que permitiu investigação rigorosa de trade-offs fundamentais entre throughput computacional (>200 pixels/s), latência determinística (<50 ms), eficiência energética (>30 GOPS/W) e precisão de classificação (>90\% Overall Accuracy). As principais contribuições incluem o desenvolvimento de metodologia reprodutível para análise comparativa de arquiteturas heterogêneas, proposição de solução técnica otimizada baseada na distribuição inteligente de tarefas entre FPGA Zynq-7020, GPU Jetson Xavier NX e ARM Cortex-A78AE, catalogação sistemática de técnicas de otimização, e estabelecimento de recomendações práticas para implementação em sistemas UAV agrícolas. As perspectivas futuras abrangem validação experimental física da arquitetura proposta, extensão para outras modalidades sensoriais, desenvolvimento de otimização dinâmica adaptativa, integração com edge computing e IoT, aplicação em domínios médicos e industriais, e criação de standards para benchmarking de sistemas hiperespectrais embarcados. O impacto esperado inclui avanços metodológicos em pesquisa acadêmica, desenvolvimento de próxima geração de sistemas UAV comerciais com maior autonomia e precisão, contribuições para sustentabilidade agrícola através de monitoramento eficiente de culturas, e estabelecimento de precedente para pesquisa interdisciplinar que integra arquiteturas computacionais, processamento de sinais e agricultura de precisão, maximizando relevância científica e aplicabilidade prática das soluções desenvolvidas.